\begin{table}[htbp]\centering
\begin{threeparttable}
\caption{Barcode fragments cut from the 904 deposited rDNA copies (amfsplit v0.3.1; primer-trimmed fragments,
primers found with up to 15\% of their length in mismatches or gaps).}
\label{tab:stage1}
\footnotesize\setlength{\tabcolsep}{3pt}
\begin{tabular}{llrrrr}
\toprule
Fragment & Used by & Usable & Missing & Open & Median (bp)\\
\midrule
SSU NS31--AM1        & MaarjAM VTs      & 876 & 28 (3$'$) & 0  & 508\\
SSU NS31--AML2       & SSU studies      & 904 & 0         & 0  & 522\\
SSU WANDA--AML2      & SSU V4           & 903 & 1 (5$'$)  & 0  & 500\\
SSU AMV4.5NF--AMDGR  & SSU V4, short    & 904 & 0         & 0  & 216\\
LSU LROR--FLR2       & LSU reference V18 & 904 & 0         & 0  & 735\\
LSU FLR3--FLR4 (D2)  & LSU D2           & 904 & 0         & 26 (3$'$) & 330\\
AML1--LSUmBr         & full amplicon    & 0   & 904 (both) & 0 & --\\
\bottomrule
\end{tabular}
\begin{tablenotes}\footnotesize
\item Usable: complete plus open-ended fragments; missing: primer absent although the copy covers its position;
open: primer site beyond the end of the copy. The 28 missing NS31--AM1 fragments are all the Paraglomerales copies,
in which AM1 does not bind; the one missing WANDA fragment is the copy dropped from the reference (PX215117.1). The
26 open D2 fragments are copies that end before the FLR4 site. The outer amplicon primers are absent because the
deposited copies are trimmed inside them. No fragment crosses an ITSxRust anchor, and none of the 904 copies is
flagged as a concatemer.
\end{tablenotes}
\end{threeparttable}
\end{table}
